\section{习题课 \#13（12月6日）}

\subsection{正定矩阵的判别}

设 $A=A^T\in\R^{n\times n}$。下列条件等价：
\begin{enumerate}
  \item 对每个 $x\neq0$，
  \begin{equation}
   x^TAx>0;
  \end{equation}
  \item $A$ 的全部特征值均为正；
  \item 存在可逆矩阵 $P$ 使
  \begin{equation}
   A=P^TP;
  \end{equation}
  \item 存在对角元为正的可逆下三角矩阵 $L$ 使
  \begin{equation}
   A=LL^T;
  \end{equation}
  \item 所有顺序主子式均为正；
  \item 所有主子式均为正。
\end{enumerate}
其中第四项是 Cholesky 分解，第五项通常称为 Sylvester 判别法。

\begin{proof}
由实对称谱定理，
\begin{equation}
 A=Q\diag(\lambda_1,\ldots,\lambda_n)Q^T.
\end{equation}
于是
\begin{equation}
 x^TAx=\sum_i\lambda_i y_i^2,
 \qquad y=Q^Tx,
\end{equation}
所以前两项等价。若所有 $\lambda_i>0$，取
\begin{equation}
 P=\diag(\sqrt{\lambda_1},\ldots,\sqrt{\lambda_n})Q^T
\end{equation}
便有 $A=P^TP$；反向由 $x^TP^TPx=\norm{Px}^2>0$ 得到。

Cholesky 分解可按 Schur 补归纳构造。写
\begin{equation}
 A=\begin{pmatrix}a&b^T\\b&C\end{pmatrix},
 \qquad a>0.
\end{equation}
则 $A$ 正定当且仅当
\begin{equation}
 C-a^{-1}bb^T
\end{equation}
正定。对该 Schur 补递归分解即可。分解中前 $k$ 阶顺序主子式等于前 $k$ 个 Cholesky 对角元平方的乘积，因此正定蕴含所有顺序主子式为正；反向可按同一递推构造 Cholesky 分解。所有主子式为正显然蕴含顺序主子式为正，而正定矩阵的任意主子矩阵仍正定，故其主子式均正。
\end{proof}

\subsection{两个正定性例题}

\begin{example}
考虑
\begin{equation}
 f(x_1,x_2,x_3)
 =99x_1^2-12x_1x_2+48x_1x_3
 +130x_2^2-60x_2x_3+71x_3^2.
\end{equation}
其对称矩阵为
\begin{equation}
 A=\begin{pmatrix}
 99&-6&24\\
 -6&130&-30\\
 24&-30&71
 \end{pmatrix}.
\end{equation}
顺序主子式为
\begin{equation}
 \Delta_1=99>0,
 \qquad
 \Delta_2=99\cdot130-36=12834>0,
\end{equation}
以及
\begin{equation}
 \Delta_3=\det A=755874>0.
\end{equation}
故 $A$ 正定，二次型 $f$ 为正定二次型。
\end{example}

\begin{example}
考虑参数矩阵
\begin{equation}
 B(t)=\begin{pmatrix}
 1&t&5\\
 t&4&3\\
 5&3&1
 \end{pmatrix}.
\end{equation}
若 $B(t)$ 正定，则前两个顺序主子式要求
\begin{equation}
 4-t^2>0,
 \qquad -2<t<2.
\end{equation}
但
\begin{equation}
 \det B(t)=-t^2+30t-105
 =-(t-15)^2+120.
\end{equation}
其为正要求
\begin{equation}
 15-2\sqrt{30}<t<15+2\sqrt{30},
\end{equation}
与 $(-2,2)$ 没有交集。因此不存在实数 $t$ 使 $B(t)$ 正定。
\end{example}

\subsection{Cholesky 分解}

\begin{theorem}[Cholesky]
若 $A=A^T\in\R^{n\times n}$ 正定，则存在唯一的下三角矩阵 $L$，其对角元全正，并满足
\begin{equation}
 A=LL^T.
\end{equation}
\end{theorem}

\begin{proof}
存在性已可由上一节的 Schur 补递推得到。为说明唯一性，若
\begin{equation}
 A=L_1L_1^T=L_2L_2^T,
\end{equation}
则
\begin{equation}
 L_2^{-1}L_1=L_2^TL_1^{-T}.
\end{equation}
左端是对角元全正的下三角矩阵，右端是上三角矩阵，所以二者必须是对角矩阵。再由
\begin{equation}
 (L_2^{-1}L_1)(L_2^{-1}L_1)^T=I
\end{equation}
及正对角元可知该对角矩阵为 $I$，故 $L_1=L_2$。
\end{proof}

\subsection{利用迹刻画正定性}

\begin{theorem}
设 $A=A^T\in\R^{n\times n}$ 可逆。则 $A$ 正定，当且仅当对每个正定矩阵 $B$ 都有
\begin{equation}
 \tr(AB)>0.
\end{equation}
\end{theorem}

\begin{proof}
若 $A$ 正定，写成
\begin{equation}
 A=P^TP,
 \qquad B=Q^TQ,
\end{equation}
其中 $P,Q$ 可逆。则
\begin{equation}
 \tr(AB)=\tr(P^TPQ^TQ)
 =\tr\bigl((QP^T)(QP^T)^T\bigr)
 =\norm{QP^T}_F^2>0.
\end{equation}

反之，若存在 $x\neq0$ 使 $x^TAx<0$，取
\begin{equation}
 B=xx^T+\varepsilon I
\end{equation}
并令 $\varepsilon>0$ 足够小，则 $B$ 正定而
\begin{equation}
 \tr(AB)=x^TAx+\varepsilon\tr A<0,
\end{equation}
矛盾。因此 $A$ 半正定。结合 $A$ 可逆可知 $A$ 正定。
\end{proof}

\subsection{实对称矩阵的相似与正交相似}

\begin{theorem}
设 $A,B\in\R^{n\times n}$ 都是实对称矩阵。则
\begin{equation}
 A\text{ 与 }B\text{ 相似}
 \quad\Longleftrightarrow\quad
 A\text{ 与 }B\text{ 正交相似}.
\end{equation}
\end{theorem}

\begin{proof}
正交相似当然蕴含相似。反之，设
\begin{equation}
 AP=PB
\end{equation}
且 $P$ 可逆。转置后有
\begin{equation}
 P^TA=BP^T.
\end{equation}
因此
\begin{equation}
 B(P^TP)=P^TAP=(P^TP)B.
\end{equation}
令
\begin{equation}
 H=(P^TP)^{1/2}.
\end{equation}
由多项式插值或函数演算，$H$ 也与 $B$ 可交换。取极分解
\begin{equation}
 P=QH,
 \qquad Q\text{ 正交}.
\end{equation}
代入 $AP=PB$ 得
\begin{equation}
 AQH=QHB=QBH.
\end{equation}
右乘 $H^{-1}$，得到 $AQ=QB$，即
\begin{equation}
 A=QBQ^T.
\end{equation}
\end{proof}

\subsection{秩一修正下的惯性变化}

对实对称矩阵 $M$，记
\begin{equation}
 n_+(M),\quad n_-(M),\quad n_0(M)
\end{equation}
分别为正、负、零特征值个数，并定义符号差
\begin{equation}
 s(M)=n_+(M)-n_-(M).
\end{equation}

\begin{theorem}
设 $A=A^T\in\R^{n\times n}$ 可逆，$\alpha\in\R^n$，并令
\begin{equation}
 B=A-\alpha\alpha^T,
 \qquad
 r=\alpha^TA^{-1}\alpha.
\end{equation}
则
\begin{equation}
 s(A)=s(B)+1-\sgn(1-r),
\end{equation}
其中约定 $\sgn(0)=0$。也就是说，
\begin{equation}
 s(A)=
 \begin{cases}
 s(B),&r<1,\\
 s(B)+1,&r=1,\\
 s(B)+2,&r>1.
 \end{cases}
\end{equation}
当 $r=1$ 时，$B$ 必奇异。
\end{theorem}

\begin{proof}
考虑对称分块矩阵
\begin{equation}
 M=\begin{pmatrix}1&\alpha^T\\\alpha&A\end{pmatrix}.
\end{equation}
以左上角 $1$ 作 Schur 补，得到
\begin{equation}
 M\sim_{\mathrm{cong}}\diag(1,B),
\end{equation}
故
\begin{equation}
 s(M)=1+s(B).
\end{equation}
另一方面，以 $A$ 作 Schur 补，得到
\begin{equation}
 M\sim_{\mathrm{cong}}\diag(1-r,A),
\end{equation}
所以
\begin{equation}
 s(M)=\sgn(1-r)+s(A).
\end{equation}
比较两式即可。
\end{proof}

\subsection{秩一全一修正矩阵的惯性}

令
\begin{equation}
 A=\lambda I_n+\ones\ones^T.
\end{equation}
向量 $\ones$ 是特征向量，对应特征值 $\lambda+n$；在 $\ones^\perp$ 上，$A$ 等于 $\lambda I$。因此特征值为
\begin{equation}
 \lambda+n\quad\text{一次},
 \qquad
 \lambda\quad\text{$n-1$ 次}.
\end{equation}
符号差为
\begin{equation}
 s(A)=
 \begin{cases}
 n,&\lambda>0,\\
 1,&\lambda=0,\\
 2-n,&-n<\lambda<0,\\
 1-n,&\lambda=-n,\\
 -n,&\lambda<-n.
 \end{cases}
\end{equation}
这也给出正、负惯性指数及奇异参数值。

\subsection{两个特殊二次型}

\begin{example}
令
\begin{equation}
 q(x)=\sum_{1\leq j<k\leq n}(-1)^{j+k}x_jx_k.
\end{equation}
取
\begin{equation}
 \alpha=(1,-1,1,-1,\ldots)^T.
\end{equation}
则 $q(x)=x^TAx$ 的矩阵满足
\begin{equation}
 A=-\frac12I+\frac12\alpha\alpha^T.
\end{equation}
因为 $\alpha^T\alpha=n$，$A$ 在 $\alpha$ 方向上的特征值为
\begin{equation}
 \frac{n-1}{2},
\end{equation}
在 $\alpha^\perp$ 上的特征值为 $-1/2$。因此 $n>1$ 时其惯性为
\begin{equation}
 (n_+,n_-,n_0)=(1,n-1,0).
\end{equation}
\end{example}

\begin{example}
令
\begin{equation}
 g(x)=\sum_{1\leq j<k\leq n}|j-k|x_jx_k.
\end{equation}
其对称矩阵 $D$ 的对角元为 $0$，非对角元为
\begin{equation}
 D_{jk}=\frac12|j-k|.
\end{equation}
依次对相邻行、列作差，可将 $D$ 合同化为一个正的一维块与 $n-1$ 个负的一维块。因此
\begin{equation}
 (n_+(D),n_-(D),n_0(D))=(1,n-1,0),
 \qquad n\geq2.
\end{equation}
原稿中的连续合同行列变换正是这一结论的直接计算证明。
\end{example}
